\section{Introduction}
\label{sec:introduction}

Language-model agents interleave generated text, external retrieval, tools,
memory, and persistent actions.  ReAct-style control loops
\citep{yao2023react} and agent frameworks such as AgentScope
\citep{gao2024agentscope} make these components easier to compose, while
benchmarks expose failures in interactive environments
\citep{liu2024agentbench}.  In medicine, however, a correct-looking final
answer is only one observable at one instant.  A response can cite a source
that does not support a particular claim; a cancelled worker can finish after
its successor; or a model, prompt, retriever, or tool update can improve an
average score while degrading an elderly or high-risk slice.  Each failure can
occur even when an evaluator assigns a plausible score to the final text.

Healthcare guidance accordingly emphasizes traceability, human factors,
versioning, monitoring, and lifecycle governance
\citep{vasey2022decide,liu2020consortai,lekadir2025futureai}.  Yet these
principles do not by themselves specify which runtime objects must be
verified, when verification occurs, or what state must be committed together.
Our central position is that high-risk agent verification should cover three
objects on three time scales:
\begin{enumerate}[leftmargin=*,nosep]
    \item a \textbf{claim}, whose admissibility depends on traceable evidence
    and deterministic risk rules;
    \item a \textbf{run}, whose public outcome depends on current ownership and
    an atomic terminal state; and
    \item an \textbf{update}, whose release depends on regression checks that
    the candidate cannot observe or approve.
\end{enumerate}

We develop this position through \system, a Web-based research prototype for
older patients and geriatric clinicians.  It uses an AgentScope runtime with
retrieval, memory, governed tools, streaming output, and an operator-run
offline evolution controller.  \system\ is an instructive case because
geriatric tasks combine multimodal evidence, polypharmacy, missing or
conflicting longitudinal facts, emergency escalation, and costly failures.
The contribution is not a new clinical model.  It is a reference verification
architecture, examined against a partial implementation, in which generated
content, concurrent execution, and candidate evolution cross explicit trust
boundaries.

This paper makes three contributions:
\begin{itemize}[leftmargin=*,nosep]
    \item We define a three-boundary model---\claimcheck,
    \runcommit, and \updategate---that separates evidential validity,
    transactional finality, and update admissibility.
    \item We report implementation status boundary by boundary: citation-marker
    auditing and bounded diagnosis rewriting; lease-plus-fencing terminal
    commits; and frozen sandbox execution, a routing-only paired runner, and
    release-control contracts.  We identify independent sealed evaluation and
    approval services as unimplemented deployment obligations.
    \item We map an adjacent deterministic audit of 40 passing policy cases to
    the claims it can and cannot support.  It is not evidence for run finality,
    update admissibility, clinical effectiveness, or general safety.
\end{itemize}

The intended value is a falsifiable systems blueprint and a set of
machine-checkable obligations for the workshop community.  No human-subject
data, clinician evaluation, diagnostic-accuracy experiment, or live clinical
deployment is reported.
